% Part: turing-machines
% Chapter: undecidability
% Section: universal-tm

\documentclass[../../../include/open-logic-section]{subfiles}

\begin{document}

\olfileid{tur}{und}{uni}
\olsection{ماشین‌های تورینگ همگانی}

در \olref[enu]{sec} بحث کردیم که چگونه هر ماشین تورینگ را می‌توان با
دنباله‌ای متناهی از اعداد صحیح توصیف کرد. این دنباله حالت‌ها، الفبا،
حالت آغازین و دستورهای ماشین تورینگ را رمزگذاری می‌کند. همچنین اشاره
کردیم که مجموعهٔ همهٔ این توصیف‌ها !!{enumerable} است. چون مجموعهٔ چنین
توصیف‌هایی !!{denumerable}، تابعی !!a{surjective} از~$\Nat$ به این توصیف‌ها وجود
دارد. برای نمونه، می‌توان چنین تابع !!a{surjective}یی را با روش زیگزاگ کانتور
به دست آورد. این تابع راهی برای برشمردن همهٔ (توصیف‌های) ماشین‌های
تورینگ به دست می‌دهد. اگر یکی از این برشمردن‌ها را تثبیت کنیم، اکنون
سخن‌گفتن از ماشین تورینگ نخست، دوم،~\dots، $e$اُم معنا دارد. این اعداد
\emph{نمایه} نامیده می‌شوند.

\begin{defn}
اگر $M$ ماشین تورینگ $e$اُم (در برشمردن تثبیت‌شدهٔ ما) باشد، می‌گوییم
$e$ \emph{نمایه‌ای} از~$M$ است. ماشین تورینگ $e$اُم را با~$M_e$
می‌نویسیم.
\end{defn}

یک ماشین ممکن است بیش از یک نمایه داشته باشد؛ مثلاً دو توصیف از~$M$
ممکن است فقط در ترتیب فهرست‌شدن دستورهایش تفاوت داشته باشند و این
توصیف‌های متفاوت نمایه‌های متفاوتی خواهند داشت.

نکتهٔ مهم این است که می‌توان توصیف‌های ماشین‌های تورینگ را چنان
برشمرد که بتوانیم توصیف~$M$ را از نمایه‌اش به‌طور مؤثر محاسبه کنیم و
از توصیف یک ماشین~$M$ نیز نمایه‌ای برای آن را به‌طور مؤثر به دست
آوریم. بنا بر تز چرچ--تورینگ، آنگاه می‌توان ماشین تورینگی یافت که
توصیف ماشین تورینگ با نمایهٔ~$e$ را بازیابی کند و توصیف متناظر را
به‌عنوان خروجی روی نوارش بنویسد. این توصیف دنباله‌ای از بلوک‌های
$\TMstroke$ خواهد بود (که اعداد صحیح مثبتِ دنبالهٔ توصیف‌کنندهٔ~$M_e$
را نمایش می‌دهند).

با این مقدمه، اکنون طبیعی است بپرسیم: کدام توابعِ نمایه‌های ماشین
تورینگ خود با ماشین‌های تورینگ محاسبه‌پذیرند؟ کدام ویژگی‌های نمایه‌های
ماشین تورینگ را ماشین‌های تورینگ می‌توانند تصمیم کنند؟ یک نمونه:
تابعی که نمایهٔ~$e$ را به تعداد حالت‌های ماشین تورینگ دارای آن نمایه
می‌فرستد با ماشین تورینگ محاسبه‌پذیر است. چنین ماشینی این‌گونه کار
می‌کند: پس از آغاز روی نواری دارای یک بلوک از~$e$ نماد~$\TMstroke$،
نخست $e$ را به توصیفش رمزگشایی می‌کند. اکنون توصیف با دنباله‌ای از
بلوک‌های~$\TMstroke$ روی نوار نمایش داده شده است. چون نخستین !!{element}
این دنباله تعداد حالت‌هاست، فقط باید همه‌چیز جز نخستین بلوک
$\TMstroke$ را پاک کند و سپس متوقف شود.

نتیجهٔ شگفت‌انگیز زیر برقرار است:

\begin{thm}\ollabel{thm:universal-tm} \emph{ماشین تورینگ همگانی‌ای}
  چون~$U$ وجود دارد که وقتی با ورودی~$\tuple{e,n}$ آغاز شود،
  \begin{enumerate}
    \item متوقف می‌شود اگر و تنها اگر $M_e$ با ورودی~$n$ متوقف شود؛ و
    \item اگر $M_e$ با خروجی~$m$ متوقف شود، $U$ نیز چنین می‌کند.
  \end{enumerate}
  پس $U$ تابع $f\colon\Nat\times\Nat\pto\Nat$ را محاسبه می‌کند که
  $f(e,n)=m$ است اگر $M_e$ پس از آغاز با ورودی~$n$ با خروجی~$m$
  متوقف شود، و در غیر این صورت تعریف‌نشده است.
\end{thm}

\begin{proof}
  ساخت واقعیِ~$U$ عملاً ناممکن است، زیرا ماشینی بسیار پیچیده خواهد
  بود. اما می‌توانیم طرح کلیِ کار آن را توصیف کنیم و سپس به تز
  چرچ--تورینگ استناد کنیم. هنگام آغاز، نوار~$U$ یک بلوک از~$e$ نماد
  $\TMstroke$ و پس از آن یک بلوک از~$n$ نماد~$\TMstroke$ دارد. نخست
  نمایهٔ~$e$ را در سمت راست ورودی~$n$ «رمزگشایی» می‌کند. نتیجه فهرستی
  از اعداد است (یعنی بلوک‌های~$\TMstroke$ جداشده با~$\TMblank$) که
  دستورهای ماشین~$M_e$ را توصیف می‌کند. سپس $U$ شمارهٔ حالت آغازین
  $M_e$ و عدد~$1$ را در سمت راست توصیف~$M_e$ روی نوار می‌نویسد. (این
  اعداد نیز یکانی و به صورت بلوک‌های~$\TMstroke$ نمایش داده می‌شوند.)
  سپس ورودی (بلوک~$n$تاییِ~$\TMstroke$) را به سمت راست کپی می‌کند، اما
  هر~$\TMstroke$ را با بلوکی از سه~$\TMstroke$ جایگزین می‌کند (به یاد
  آورید شمارهٔ نماد~$\TMstroke$ عدد~$3$ است؛ $1$ شمارهٔ~$\TMendtape$
  و~$2$ شمارهٔ~$\TMblank$ است). در انتهای چپ این دنبالهٔ بلوک‌ها (که
  روی نوار~$U$ با نمادهای~$\TMblank$ جدا شده‌اند) یک~$\TMstroke$، یعنی
  رمز~$\TMendtape$، می‌نویسد.

  اکنون روی نوار~$U$ این‌ها قرار دارند: نمایهٔ~$e$، عدد~$n$، شمارهٔ
  رمز حالت آغازین («حالت فعلی»)، عدد~$1$ برای مکان آغازین سر («مکان
  فعلی سر») و محتوای آغازین «نوار» (دنباله‌ای از بلوک‌های~$\TMstroke$
  که شماره‌رمزهای نمادهای~$M_e$، یعنی «نمادها»، را نمایش می‌دهند و با
  $\TMblank$ از هم جدا شده‌اند).

  سپس با انجام کارهای زیر آنچه را $M_e$ پس از آغاز با ورودی~$n$ انجام
  می‌دهد شبیه‌سازی می‌کند:
  \begin{enumerate}
    \item عدد~$k$ مربوط به «مکان فعلی سر» را بیابد (در آغاز این عدد
      $1$ است)؛
    \item به بلوک $k$اُم در «نوار» برود و ببیند «نماد» آنجا چیست؛
    \item\ollabel{find-inst}%
      دستور متناظر با «حالت» و «نماد» فعلی را بیابد؛
    \item به بلوک $k$اُمِ «نوار» بازگردد و «نماد» آنجا را با شمارهٔ
      رمز نمادی که $M_e$ می‌نوشت جایگزین کند؛
    \item سر را به جایی ببرد که «حالت» فعلی ثبت شده است و عدد آنجا را
      با شمارهٔ حالت جدید جایگزین کند؛
    \item به جایی برود که «مکان نوار» ثبت شده است و یک~$\TMstroke$
      پاک یا اضافه کند (به‌ترتیب اگر دستور حرکت به چپ یا راست باشد)؛
    \item تکرار کند.\footnote{در اینجا از بعضی دشواری‌های ظریف چشم
      می‌پوشیم. برای نمونه، ممکن است $U$ هنگام افزایش شمارنده‌ای که
      «مکان فعلی سر» را در آن نگه می‌دارد به فضای بیشتری نیاز داشته
      باشد؛ در آن صورت باید تمام «نوار» را به راست جابه‌جا کند.}
  \end{enumerate}
  اگر $M_e$ پس از آغاز با ورودی~$n$ هرگز متوقف نشود، $U$ نیز هرگز
  متوقف نمی‌شود و خروجی‌اش تعریف‌نشده است.

  اگر در گام~\olref{find-inst} معلوم شود توصیف~$M_e$ هیچ دستوری برای
  زوج «حالت»/«نماد» فعلی ندارد، $M_e$ متوقف می‌شد. اگر چنین شود، $U$
  بخش نوارش را که در سمت چپ «نوار» قرار دارد پاک می‌کند. نخست بلوک
  یک‌تاییِ~$\TMstroke$ را که رمز~$\TMendtape$ است بررسی می‌کند و از آن می‌گذرد؛
  اگر بلوک نخست این رمز را نداشته باشد، ورودی رمزگذاری‌شده را رد
  می‌کند. سپس تا رسیدن به نخستین بلوک دوتاییِ~$\TMstroke$، که
  رمز~$\TMblank$ است، هر بلوک سه‌تاییِ~$\TMstroke$ را به یک
  $\TMstroke$ در انتهای چپ نوار خودش تبدیل می‌کند. با دیدن نخستین رمز
  $\TMblank$، آن بلوک و همهٔ باقی‌ماندهٔ رمزگذاری «نوار» را پاک
  می‌کند. پس از پایان این کار، نوار~$U$ شامل یک بلوک واحدِ
  $\TMstroke$ به طول~$m$ است.

  اگر $U$ پس از رمز~$\TMendtape$ و پیش از نخستین رمز~$\TMblank$ با
  رمزی جز بلوک سه‌تاییِ~$\TMstroke$ روبه‌رو شود، بی‌درنگ آن رمزگذاری
  را رد و متوقف می‌شود. چون در این حالت نوار~$U$ شامل یک بلوک واحدِ
  $\TMstroke$ نیست، خروجی آن عدد طبیعی نیست؛ یعنی $f(e,n)$ در این
  حالت تعریف‌نشده است.
\end{proof}

\end{document}
